\chapter{Introduction}
\label{ch:introduction}

Algorithms run an increasing share of bidding in online auctions. Programmatic advertising, sponsored search, electricity markets, transport pricing, and public procurement all now rely on autonomous agents that submit bids on behalf of human principals. When these agents learn from repeated play rather than compute equilibrium strategies, the classical auction theorems that govern first-price and second-price formats no longer apply directly, and the policy recommendations that flow from those theorems may not apply either. Competition authorities have flagged this gap in repeated roundtables \citep{OECD2017, EzrachiStucke2017, OECD2023}, but the empirical basis for any specific regulatory response remains thin.

This thesis asks a narrow economic question. Across a common experimental framework, how do auction format, competitive thickness, learning parameters, and budget constraints rank in their association with platform revenue, and do these rankings generalise across algorithm classes? Answering the question requires running algorithmic bidders under controlled conditions. Empirical data on real-world algorithmic auctions are scarce because most auctions are proprietary, and theoretical analysis faces a parallel barrier, since modelling the kinds of bidding algorithms deployed in practice, which use large tables of learned values or continually adjust shadow prices to keep spend within a budget, is analytically intractable beyond small special cases. I therefore build a computational laboratory based on factorial experimental design and large-scale Monte Carlo simulation, a style of statistical experiment in which the inputs (auction rules, algorithm parameters, market conditions) are systematically varied and the outputs (revenue, efficiency) are regressed on the inputs. The laboratory supplies the controlled variation that theory cannot tractably derive and that proprietary markets do not make available.

\section{The Two-Literature Tension}
\label{sec:intro_tension}

The thesis is motivated by a tension between two bodies of work that reach opposite conclusions about auction format. On one side, reinforcement-learning studies of algorithmic collusion conclude that first-price auctions are more vulnerable to supra-competitive bid suppression than second-price auctions \citep{BanchioSkrzypacz2022, Calvano2020, Klein2021}. On the other, the autobidding and pacing literature concludes that first-price auctions achieve tighter welfare guarantees and often generate higher realised revenue under budget constraints \citep{Aggarwal2019, Balseiro2019Learning, Conitzer2022FPPE, Gaitonde2023}. These findings sit directly in opposition. The former suggests that platforms should prefer second-price formats to protect sellers from algorithmic coordination, while the latter suggests that platforms should prefer first-price formats to achieve welfare guarantees under budget constraints. The two literatures reach opposite conclusions because they study different bidding technologies, not because one is wrong. This thesis tests both claims within a single experimental framework and asks whether each holds in the setting its original authors studied and whether either generalises to the other.

\section{What This Thesis Does}
\label{sec:intro_what}

I run six sub-experiments, organised into three families indexed by the algorithm governing bidder behaviour. Experiments~1a and~1b use tabular Q-learning, a reinforcement-learning rule that stores a table of estimated action values and updates each entry from observed rewards, under constant and affiliated valuations respectively.\footnote{Reinforcement learning is the branch of machine learning in which an agent interacts with an environment over time, receives numerical rewards, and adjusts its behaviour to increase those rewards. Q-learning is the oldest and simplest such rule.} Experiments~2a and~2b replace Q-learning with contextual bandits, a simpler class of learning algorithms that picks an action from a set of alternatives based on an observed context vector and a running estimate of expected payoffs. The two contextual-bandit variants I study are LinUCB, which adds an optimism bonus to an estimate of expected reward, and Thompson Sampling, which draws a posterior sample of the reward parameters and picks the action that is best under the sample. Experiments~3a and~3b replace unconstrained learners with budget-constrained pacing agents, rules that shade bids proportionally to how quickly the agent is spending its budget so that the budget lasts the full horizon. The two pacing variants are multiplicative dual pacing, which adjusts a Lagrangian shadow price on the remaining budget, and a proportional-integral controller, a feedback rule borrowed from control engineering that reacts to accumulated spending error. Every experiment uses a factorial or fractional factorial design that varies auction format, the number of bidders, and between four and eight additional factors that are specific to the algorithm family. Across the six sub-experiments I report ranked effects for platform revenue, price volatility, and allocative efficiency on more than three thousand independent simulation runs.

I treat each simulation as a black-box observation. I vary inputs, measure outcomes, and report which factors associate with larger or smaller revenue shortfalls, in what direction, at what magnitude, and how these associations differ across algorithm classes. I do not attempt to explain the mechanisms by which algorithms produce observed outcomes. I do not investigate whether agents develop reward-punishment strategies, whether observed bid suppression reflects tacit coordination or merely suboptimal convergence, or whether budget-rational bid shading is involved. An algorithm that fails to explore sufficiently, one that rationally manages a finite budget, and one that learns to punish deviators can all produce identical revenue shortfalls in outcome data. I use the term bid suppression throughout to describe an observed revenue gap without presupposing its cause.

\section{Headline Findings}
\label{sec:intro_findings}

Two cross-cutting patterns emerge consistently. First, structural market parameters dominate algorithmic design choices. In unconstrained settings, the number of bidders is the strongest predictor of revenue, consistent with the direction of \citet{BulowKlemperer1996}. Under budget constraints, the budget multiplier replaces the number of bidders as the dominant factor. Algorithmic tuning parameters, including learning rates, discount factors, regularisation, memory decay, exploration intensity, and controller aggressiveness, are collectively negligible in every experiment. Second, the effect of auction format on revenue reverses sign across algorithm classes. Under Q-learning with constant valuations, the two formats generate nearly identical revenue. Under affiliated-signal Q-learning and both contextual-bandit algorithms, second-price auctions generate higher revenue. Under both pacing variants, first-price auctions generate higher revenue. The format reversal holds under all thirteen robustness checks reported in Appendix~\ref{ch:robustness}. \Cref{fig:intro_forest} summarises both patterns across the six sub-experiments. The shift in colour and shape marks the reversal from learning experiments to pacing experiments.

\begin{figure}[H]
\centering
\includegraphics[width=0.9\textwidth]{intro_forest}
\caption{Standardised marginal effects of competition and auction format on platform revenue across the six sub-experiments. Each point is an ordinary-least-squares regression coefficient normalised by the experiment grand mean, and bars show 95 percent confidence intervals. The colour of each point encodes the total-order Sobol index $S_T$, a variance-decomposition measure that gives the share of outcome variance attributable to the factor including all of its interactions. Circles mark learning experiments and diamonds mark pacing experiments. The reversal of the auction-format sign between the two groups is the central finding of the thesis.}
\label{fig:intro_forest}
\end{figure}

Because the effect of auction format depends on the bidding technology, no single format is universally superior when bidders are algorithms. Format recommendations derived from one algorithm class do not generalise to another. This is the central methodological and policy message of the thesis.

\section{Roadmap}
\label{sec:intro_roadmap}

Chapter~\ref{ch:literature} reviews the four bodies of work that motivate the thesis and identifies the dozen anchor papers whose predictions I engage. Chapter~\ref{ch:auctions_and_algorithms} sets up the auction environment and describes the three algorithm families in economically interpretable terms, with formal update rules for reproducibility. Chapter~\ref{ch:design_and_inference} presents the factorial experimental design and the statistical inference procedure I apply uniformly to all six sub-experiments. Chapters~\ref{ch:qlearning} through~\ref{ch:pacing} report the results, one chapter per algorithm family. Chapter~\ref{ch:discussion} synthesises findings across the six sub-experiments, draws policy implications, and outlines limitations and future work. Appendices~\ref{ch:equilibria}, \ref{ch:robustness}, and~\ref{ch:ranked_tables} collect the equilibrium benchmarks, the full set of robustness diagnostics, and the ranked-effects tables for all six sub-experiments.
